\chapter{Los números hasta 100}\label{ch:g1:tens}

Contar un montón grande de uno en uno no se acaba nunca. El gran truco:
agrupar los objetos en paquetes de \emph{diez} y contar los paquetes. Así
están construidos los números hasta el $100$ --- y por eso se
escriben con dos cifras.

\section{Paquetes de diez}

\begin{definition}[Decenas y unidades]\label{def:g1:tens:def}
Una \emph{decena} es un paquete de diez unidades. El número $34$ significa $3$
decenas y $4$ unidades:
\[
34 = 10 + 10 + 10 + 4 .
\]
La cifra de la izquierda cuenta los paquetes; la de la derecha, las unidades sueltas.
\end{definition}

\begin{omfigure}
\begin{tikzpicture}[scale=0.75]
  % tres paquetes de diez palillos, cada uno atado con una goma
  \foreach \b in {0,1,2} {
    \begin{scope}[xshift=\b*2.4cm]
    \foreach \i in {0,...,9} \draw[omDef, very thick]
      ({0.16*\i},0) -- ({0.16*\i},1.7);
    \draw[omExo, thick, rounded corners=3pt]
      (-0.22,0.62) rectangle (1.66,1.05);
    \end{scope}
  }
  % cuatro palillos sueltos
  \foreach \i in {0,...,3} \draw[omThm, very thick]
    (7.7+0.4*\i,0) -- (7.7+0.4*\i,1.7);
  \node[right, font=\small] at (9.7,0.85)
    {$3$ paquetes y $4$ palillos: $34$};
\end{tikzpicture}

{\small Treinta y cuatro palillos: tres paquetes de diez (atados cada uno
con una goma) y cuatro palillos sueltos. ¡No hace falta contar hasta $34$ de uno en uno!}
\end{omfigure}

\begin{example}[Leer los números]\label{ex:g1:tens:reading}
$20$ son dos decenas (veinte), $30$ tres decenas (treinta)\,\dots\ hasta
$90$ (noventa) y $100$ (cien: ¡diez decenas!). Entre medias:
$47$ se lee «cuarenta y siete» --- cuarenta por las $4$ decenas y siete por
las unidades.
\end{example}

\begin{omfigure}
\begin{tikzpicture}[scale=0.9]
  \draw (0,0) grid[xstep=2.2, ystep=0.8] (4.4,1.6);
  \node[font=\small] at (1.1,1.2) {decenas};
  \node[font=\small] at (3.3,1.2) {unidades};
  \node[font=\large, omDef] at (1.1,0.4) {$3$};
  \node[font=\large, omThm] at (3.3,0.4) {$4$};
\end{tikzpicture}

{\small La tabla de lugares del $34$: el $3$ ocupa el lugar de las decenas y el
$4$ el de las unidades. ¡Esta tablita gana una columna nueva cada
año!}
\end{omfigure}

\begin{example}[Las mismas cifras, números distintos]\label{ex:g1:tens:place}
$27$ y $72$ usan las mismas cifras, pero son muy distintos: $27$ son
$2$ decenas y $7$ unidades; $72$ son $7$ decenas y $2$ unidades --- ¡mucho
mayor! El \emph{lugar} de una cifra decide lo que cuenta.
\end{example}

\section{La tabla del cien}

\begin{omfigure}
\begin{tikzpicture}[scale=0.52]
  \foreach \r in {0,...,9} {
    \foreach \c in {0,...,9} {
      \pgfmathtruncatemacro{\n}{\r*10+\c+1}
      \draw[gray!50] (\c,-\r) +(-0.5,-0.5) rectangle +(0.5,0.5);
      \node[font=\footnotesize] at (\c,-\r) {\n};
    }
  }
  \draw[omThm, very thick] (3,-4) +(-0.5,-0.5) rectangle +(0.5,0.5);
  \draw[omDef, very thick] (3,-5) +(-0.5,-0.5) rectangle +(0.5,0.5);
\end{tikzpicture}

{\small La tabla del cien: cada fila es una decena. Un paso a la derecha es
$+1$; un paso hacia \emph{abajo} es $+10$ (del $44$ rojo, bajando en
línea recta, se llega al $54$ azul).}
\end{omfigure}

\begin{method}[Moverse por la tabla]\label{met:g1:tens:grid}
\begin{enumerate}
  \item $+1$ / $-1$: un paso a la derecha / a la izquierda;
  \item $+10$ / $-10$: un paso hacia abajo / hacia arriba --- solo cambia la
        cifra de las decenas: $37 + 10 = 47$;
  \item contar de diez en diez: $10, 20, 30, \dots$ --- basta con seguir la
        última columna de la tabla hacia abajo.
\end{enumerate}
\end{method}

\begin{example}[Contar de 2 en 2 y de 5 en 5]\label{ex:g1:tens:skip}
De $2$ en $2$: $2, 4, 6, 8, 10, 12, \dots$ (los números pares). De $5$ en $5$:
$5, 10, 15, 20, 25, \dots$ --- terminan en $5$ o en $0$ y dibujan un
patrón bonito en la tabla. Contar a saltos prepara la multiplicación
(\cref{ch:g2:mult}).
\end{example}

\section{Comparar números mayores}

\begin{method}[Comparar con las decenas]\label{met:g1:tens:compare}
Gana quien tiene más decenas: $61 > 58$ porque $6$ decenas ganan a $5$ decenas ---
¡aunque $8$ sea más que $1$! Solo cuando las decenas son iguales deciden
las unidades: $43 < 47$.
\end{method}

\begin{example}\label{ex:g1:tens:compare}
Ordena $35$, $53$ y $29$: primero las decenas ($2$, $3$, $5$):
\[
29 < 35 < 53 .
\]
\end{example}

\section{Ejercicios}

\begin{exercise}[$\star$]\label{exo:g1:tens:1}
¿Cuántas decenas y cuántas unidades hay en: $26$; \ $40$; \ $73$; \ $99$?
\end{exercise}

\begin{exercise}[$\star$]\label{exo:g1:tens:2}
Escribe el número: $5$ decenas y $2$ unidades; \ $8$ decenas; \ $1$ decena y
$9$ unidades; \ $10$ decenas.
\end{exercise}

\begin{exercise}[$\star$]\label{exo:g1:tens:3}
Escribe con cifras: sesenta y tres; \ cuarenta; \ ochenta y cinco; \ cien.
\end{exercise}

\begin{exercise}[$\star$]\label{exo:g1:tens:4}
Calcula moviéndote por la tabla del cien: $23 + 10$; \quad $67 + 10$;
\quad $45 - 10$; \quad $80 - 10$.
\end{exercise}

\begin{exercise}[$\star$]\label{exo:g1:tens:5}
Continúa: $10, 20, 30, ?, ?, ?$. \quad
$5, 10, 15, ?, ?, ?$. \quad $12, 14, 16, ?, ?, ?$.
\end{exercise}

\begin{exercise}[$\star$]\label{exo:g1:tens:6}
Copia y completa con $<$ o $>$:
\[
34 \;?\; 43, \qquad
58 \;?\; 61, \qquad
70 \;?\; 69, \qquad
99 \;?\; 100 .
\]
\end{exercise}

\begin{exercise}[$\star$]\label{exo:g1:tens:7}
Ordena de menor a mayor: $46$; \ $19$; \ $64$; \ $91$; \
$41$.
\end{exercise}

\begin{exercise}[$\star$]\label{exo:g1:tens:8}
Un jardinero agrupa $52$ flores en ramos de diez. ¿Cuántos ramos
completos hace y cuántas flores sueltas le quedan?
\end{exercise}

\begin{exercise}[$\star\star$]\label{exo:g1:tens:9}
En la tabla del cien, sal del $36$; baja dos pasos y da un paso a la
izquierda. ¿Dónde caes? ¿Qué cuenta acabas de hacer?
\end{exercise}
